\documentclass[11pt]{article}
\usepackage[margin=1in]{geometry}
\usepackage{amsmath,amssymb,amsthm,mathtools,bm}
\usepackage{enumitem}
\usepackage{booktabs}
\usepackage{hyperref}
\usepackage{xcolor}

\newtheorem{theorem}{Theorem}
\newtheorem{lemma}[theorem]{Lemma}
\newtheorem{proposition}[theorem]{Proposition}
\newtheorem{assumption}[theorem]{Assumption}
\newtheorem{remark}[theorem]{Remark}
\newtheorem{definition}[theorem]{Definition}

\newcommand{\R}{\mathbb{R}}
\newcommand{\E}{\mathbb{E}}
\newcommand{\Pp}{\mathbb{P}}
\newcommand{\norm}[1]{\left\lVert #1\right\rVert}
\newcommand{\ip}[2]{\left\langle #1,#2\right\rangle}
\newcommand{\wtO}{\widetilde{O}}
\newcommand{\one}{\mathbf{1}}
\newcommand{\TV}{\mathrm{TV}}
\newcommand{\KL}{\mathrm{KL}}

\title{Simulation results for diffusion extension}
\author{}
\date{}

\begin{document}
\maketitle

\providecommand{\DiffusionFigRoot}{.}

\section{Numerical experiments for the private diffusion sampler}
\label{sec:diffusion-experiments}

We study separately the effects of the privacy budget, target smoothness, and
reverse-process discretization.  This order mirrors the error decomposition in
the theory: the private score and histogram releases determine a private clean
density, early stopping smooths that density, and the reverse diffusion
algorithm approximates the early-stopped Ornstein--Uhlenbeck (OU) law.  The
experiments therefore report a different metric for each stage rather than
combining density, sampling, and stability criteria in a single plot.

\subsection{Theory and experimental protocol}
\label{sec:diffusion-experimental-protocol}

For dimension $d=1$, the total-variation upper bound established earlier has
the polynomial part
\begin{equation}
  \delta_n(\alpha,\gamma)
  = \max\left\{
      n^{-(\gamma+1)/(2\gamma+3)},
      (n\alpha)^{-(\gamma+1)/(\gamma+2)}
    \right\}.
  \label{eq:experiment-delta}
\end{equation}
The same quantity controls the accuracy of the discretized private diffusion
sampler up to logarithmic factors when
$T=\log n$, $\tau=n^{-1}$, and $\kappa=\delta_n^2$; the corresponding grid uses
$N=\widetilde O(\delta_n^{-2})$ score evaluations.  The theory gives a
class-wise upper bound.  It does not predict exact equality of slopes for a
single finite target, nor does it order the finite-sample error levels of
different target distributions.

We use the finite multiscale family
\begin{equation}
  \rho_{\gamma,K}(x) \propto
  \exp\left\{-\frac{x^2}{2}
  - a\sum_{k=0}^{K}2^{-k(1+\gamma)}\cos(2^k x)\right\},
  \qquad K=40,\quad a=0.25.
  \label{eq:multiscale-target}
\end{equation}
The score parameters are sampled using projected Langevin Monte Carlo (LMC).  The total privacy budget is divided equally
between the score release and the histogram release.  The released histogram
and score determine a cellwise affine density $\widetilde\rho_D$.  We evaluate
its OU convolution using truncated Gaussian moments and generate samples with
the exponential reverse integrator on the grid described above.  The
implementation uses $100{,}000$ projected-LMC steps, $20{,}000$ generated
particles, and a $50{,}000$-point independent reference sample per setting.
The OU score is evaluated using a $3{,}201$-point interpolation whose error was
checked against direct truncated-moment evaluation.

All panels use
\[
  n\in\{1000,2000,4000,8000,16000\}
\]
and $30$ independent training datasets per displayed setting.  Points show
replication means and pointwise $95\%$ Student-$t$ intervals.  Empirical
sample-size slopes are ordinary least-squares fits of
$\log(\text{mean error})$ on $\log n$ using all five sample sizes.  Their
$95\%$ intervals are based on $2{,}000$ bootstrap resamples of the independent
datasets.  These intervals quantify simulation variability in a descriptive
slope; they are not confidence intervals for the theoretical exponent.

We distinguish four evaluation quantities.  First,
$\operatorname{TV}(\widetilde\rho_D,\rho_0)$ measures private density
estimation.  Second,
$\operatorname{TV}(P_\tau\widetilde\rho_D,\rho_0)$ measures the
early-stopped OU error.  Third, generated binned TV is a finite-resolution
two-sample diagnostic for the final particles.  Finally, reverse-to-OU binned
TV compares reverse samples with independent samples from
$P_\tau\widetilde\rho_D$ and isolates numerical sampling error at the chosen
resolution.

\subsection{Privacy-budget panel}
\label{sec:alpha-panel}

We first fix $\gamma=0.5$ and vary
$\alpha\in\{0.5,1,2,4,8\}$.  Within each $(n,r)$ pair, where $r$ indexes the
replication, the same training dataset and the same independent test sample
are reused at all privacy budgets.  Projected-LMC, histogram, and reverse
sampling randomness are independent across privacy levels.  This pairing
removes between-dataset variation from comparisons across $\alpha$ without
coupling the private mechanisms themselves.

For $\gamma=0.5$, Equation~\eqref{eq:experiment-delta} becomes
\begin{equation}
  \delta_n(\alpha,0.5)
  =\max\{n^{-3/8},(n\alpha)^{-3/5}\}.
  \label{eq:alpha-panel-rate}
\end{equation}
The crossover is $\alpha_\star=n^{-3/8}$, which is below every privacy budget
in the displayed range.  Thus the relevant asymptotic sample-size reference is
$-3/8$, whereas the leading polynomial bound has zero $\alpha$ exponent in
this regime.  This does not imply constant finite-sample error in $\alpha$:
hidden constants, logarithmic factors, and lower-order privacy terms remain.

Figure~\ref{fig:alpha-density-vs-alpha} shows the direct effect of relaxing
privacy.  At $n=1000$, the mean private-density TV decreases from $0.132$ at
$\alpha=0.5$ to $0.065$ at $\alpha=8$; at $n=16000$, it decreases from
$0.0289$ to $0.0206$.  The weaker dependence at the largest sample size is
consistent with the non-private term becoming dominant.  Quantitatively, the
empirical $\alpha$ slope of density TV changes from $-0.252$ at $n=1000$ to
$-0.119$ at $n=16000$.  We interpret this as finite-sample attenuation of the
privacy contribution, rather than as evidence for an exact power law in
$\alpha$.

\begin{figure}[t]
  \centering
  \includegraphics[width=0.82\linewidth]{\DiffusionFigRoot/alpha_panel/figures/density_tv_vs_alpha.pdf}
  \caption{Private-density total variation versus the privacy budget.  Each
  curve fixes the sample size.  Training datasets are paired across privacy
  budgets within a replication; error bars are pointwise $95\%$ Student-$t$
  intervals over $30$ datasets.  The horizontal axis is logarithmic.}
  \label{fig:alpha-density-vs-alpha}
\end{figure}

The sample-size comparison is clearest after the early-stopping OU step.
Figure~\ref{fig:alpha-ou-vs-n} and Table~\ref{tab:alpha-ou-slopes} show that the
OU-TV slopes approach the reference $-3/8$ as privacy is relaxed.  At
$\alpha=2,4,8$, the estimates are $-0.391$, $-0.367$, and $-0.370$,
respectively.  At $\alpha=0.5$ and $1$, the curves are steeper over this finite
range.  A steeper fitted slope does not improve the minimax upper bound; it
indicates that these settings are still moving between privacy-dominated and
sample-dominated finite-sample behavior.  The private-density slopes are also
steeper than $-3/8$ (from $-0.544$ to $-0.409$ across the five privacy
budgets), which reinforces the need to describe the theoretical exponent as a
reference rather than a fitted equality.

\begin{figure}[t]
  \centering
  \includegraphics[width=0.82\linewidth]{\DiffusionFigRoot/alpha_panel/figures/ou_tv_vs_n.pdf}
  \caption{Early-stopped OU total variation versus sample size in the privacy
  panel.  Both axes are logarithmic.  The high-privacy-budget curves have
  empirical slopes close to the theoretical sample-size reference $-3/8$.}
  \label{fig:alpha-ou-vs-n}
\end{figure}

\begin{table}[t]
  \centering
  \caption{Empirical sample-size slopes for early-stopped OU TV in the privacy
  panel.  Intervals are percentile bootstrap intervals over paired datasets.}
  \label{tab:alpha-ou-slopes}
  \begin{tabular}{@{}cccc@{}}
    \toprule
    $\alpha$ & Empirical slope & $95\%$ interval & Theory reference \\
    \midrule
    $0.5$ & $-0.512$ & $[-0.537,-0.488]$ & $-0.375$ \\
    $1$   & $-0.441$ & $[-0.461,-0.418]$ & $-0.375$ \\
    $2$   & $-0.391$ & $[-0.410,-0.371]$ & $-0.375$ \\
    $4$   & $-0.367$ & $[-0.385,-0.347]$ & $-0.375$ \\
    $8$   & $-0.370$ & $[-0.390,-0.350]$ & $-0.375$ \\
    \bottomrule
  \end{tabular}
\end{table}

Generated binned TV follows the same qualitative privacy ordering, but its
large-$n$ curves flatten as the error approaches the finite-particle and
binning resolution.  For example, at $n=16000$ the means are $0.0222$,
$0.0193$, and $0.0194$ for $\alpha=0.5,2,8$.  We therefore use density TV and
OU TV for the rate comparison and treat generated binned TV and Wasserstein
distance as supporting diagnostics.  No total-variation exponent is assigned
to Wasserstein distance.

\subsection{Smoothness panel}
\label{sec:gamma-panel}

We next fix $\alpha=2$ and vary
$\gamma\in\{0.2,0.4,0.6,0.8,1\}$.  The values $K=40$ and $a=0.25$ remain
fixed.  Unlike the privacy panel, changing $\gamma$ changes the target in
Equation~\eqref{eq:multiscale-target}; datasets are therefore redrawn
independently for every $(n,\gamma,r)$ setting.  Consequently, the vertical
ordering of the five curves is not a controlled monotonicity comparison.  In
particular, the theory does not require the $\gamma=1$ curve to lie below all
others at every finite $n$.  The relevant comparison is whether the
sample-size exponent changes with $\gamma$ as predicted.

At fixed $\alpha=2$, the sample-size reference is
\begin{equation}
  s_{\rm th}(\gamma)=-\frac{\gamma+1}{2\gamma+3}.
  \label{eq:gamma-slope-reference}
\end{equation}
Figure~\ref{fig:gamma-ou-vs-n} shows a systematic reduction in OU TV as $n$
increases for every target.  More directly,
Figure~\ref{fig:gamma-ou-slopes} and Table~\ref{tab:gamma-ou-slopes} compare the
empirical and theoretical exponents.  The theoretical reference lies within
the bootstrap interval at all five smoothness levels.  The observed slopes
also become more negative with $\gamma$, from $-0.355$ at $\gamma=0.2$ to
$-0.416$ at $\gamma=1$, closely tracking the theoretical change from $-0.353$
to $-0.400$.

\begin{figure}[t]
  \centering
  \includegraphics[width=0.82\linewidth]{\DiffusionFigRoot/gamma_panel/figures/ou_tv_vs_n.pdf}
  \caption{Early-stopped OU total variation versus sample size for the
  smoothness panel.  Changing $\gamma$ changes the target distribution, so
  curve heights should not be interpreted as a monotonicity claim in
  smoothness.  Both axes are logarithmic.}
  \label{fig:gamma-ou-vs-n}
\end{figure}

\begin{figure}[t]
  \centering
  \includegraphics[width=0.82\linewidth]{\DiffusionFigRoot/gamma_panel/figures/ou_tv_slopes.pdf}
  \caption{Empirical sample-size slopes of early-stopped OU TV and the
  theoretical reference $-(\gamma+1)/(2\gamma+3)$.  Bars are $95\%$
  dataset-bootstrap intervals.  This slope comparison, rather than the
  vertical ordering in Figure~\ref{fig:gamma-ou-vs-n}, is the primary
  smoothness result.}
  \label{fig:gamma-ou-slopes}
\end{figure}

\begin{table}[t]
  \centering
  \caption{Empirical and theoretical sample-size slopes in the smoothness
  panel, using early-stopped OU TV.}
  \label{tab:gamma-ou-slopes}
  \begin{tabular}{@{}cccc@{}}
    \toprule
    $\gamma$ & Empirical slope & $95\%$ interval & Theory reference \\
    \midrule
    $0.2$ & $-0.355$ & $[-0.370,-0.338]$ & $-0.353$ \\
    $0.4$ & $-0.379$ & $[-0.400,-0.359]$ & $-0.368$ \\
    $0.6$ & $-0.400$ & $[-0.419,-0.380]$ & $-0.381$ \\
    $0.8$ & $-0.410$ & $[-0.430,-0.389]$ & $-0.391$ \\
    $1.0$ & $-0.416$ & $[-0.437,-0.394]$ & $-0.400$ \\
    \bottomrule
  \end{tabular}
\end{table}

The density-TV slopes are uniformly steeper than their theoretical references,
whereas the generated binned-TV slopes become flatter at the smoother targets.
The latter effect is expected once generated binned TV approaches the roughly
$0.013$ two-sample resolution floor.  Early-stopped OU TV lies between these
two numerical effects and gives the cleanest comparison with
Equation~\eqref{eq:gamma-slope-reference}; we therefore designate it as the
primary metric for the smoothness-rate claim.

\subsection{Reverse-diffusion discretization}
\label{sec:reverse-discretization}

The preceding panels are informative about statistical error only if the
reverse sampler is sufficiently accurate.  Across the full panels, the mean
reverse-to-OU binned TV is approximately $0.013$, essentially the same as the
binned TV obtained by comparing two independent OU samples of equal size.
The paired reverse-minus-IID difference is close to zero in nearly every
$(n,\alpha)$ cell.  Thus, at the resolution of the experiment, the reverse
integrator is not the dominant error source.

We also perform a pre-specified discretization experiment at $n=4000$ and
$\alpha=1$ using three fixed private densities.  We vary the grid multiplier
so that the number of reverse steps ranges from $66$ to $31{,}366$, and run
$10$ independent reverse samplers with $20{,}000$ particles for each density
and grid.  Figure~\ref{fig:discretization} shows a clear reduction from about
$0.025$--$0.027$ at $N=66$ to the $0.013$--$0.014$ IID OU floor by roughly
$N=250$.  Further refinement does not yield a visible improvement because the
finite-particle/binning floor dominates.

\begin{figure}[t]
  \centering
  \includegraphics[width=\linewidth]{\DiffusionFigRoot/alpha_panel/figures/discretization.pdf}
  \caption{Reverse-process discretization sensitivity for three pre-selected
  private densities.  The plotted metric is reverse-to-OU binned TV; the IID
  OU line estimates the finite-particle/binning floor using the same sample
  size.  The horizontal axis is logarithmic.}
  \label{fig:discretization}
\end{figure}

The coarse-grid empirical slopes are between $-0.23$ and $-0.26$, whereas the
nominal discretization upper-bound reference is $-1/2$.  The all-grid slopes
are only about $-0.07$ to $-0.09$ because most points already lie on the
empirical floor.  We therefore use this experiment to establish numerical
adequacy, rather than to claim an empirical verification of the
discretization exponent.

\subsection{Single-record stability diagnostic}
\label{sec:replacement-stability}

We recommend placing this subsection and its figure in the appendix because
it studies a mechanism-level phenomenon rather than the main accuracy rate.
At $\gamma=0.5$, we compare the response to a one-record replacement for a
non-private constrained score-matching ERM, the mode of the private target,
and a projected-LMC release.  The basis, bandwidth, truncation radius, and
constraints are frozen across methods.  We use
$n\in\{1000,4000,16000\}$, $\alpha\in\{0.5,2,8\}$, and $10$ independent base
datasets; inference is performed at the dataset level.

For the two neighboring datasets, the private LMC chains use the same Gaussian
innovations.  Under this synchronous coupling, the mean output-function RMS
shift is below the corresponding non-private shift in all displayed cells.
For typical replacements, the ratio ranges from $0.089$ to $0.637$; for the
pre-specified stress replacements it ranges from $0.022$ to $0.154$.  The
ratio generally increases with $\alpha$, as expected when privacy noise is
reduced.  The deterministic private mode shows a more modest stability gain,
especially for typical replacements.

\begin{figure}[t]
  \centering
  \includegraphics[width=\linewidth]{\DiffusionFigRoot/stability/figures/typical_private_coupled_function_rms.pdf}
  \caption{Output-function RMS change after a typical one-record replacement.
  Private chains use synchronous Gaussian innovations for the neighboring
  datasets.  Coefficient-space diagnostics are reported separately in the
  supplement.}
  \label{fig:replacement-stability}
\end{figure}

This experiment is deliberately described as a diagnostic.  Differential
privacy controls output distributions for neighboring datasets; it does not
require every synchronously coupled realization to move less than a
non-private estimator.  Moreover, two independent LMC releases from the same
dataset exhibit a mechanism-randomness floor that is larger than the coupled
input response in many settings.  Figure~\ref{fig:replacement-stability}
therefore illustrates reduced input sensitivity under one controlled coupling
and does not constitute a separate privacy guarantee.

\subsection{Summary and limitations}
\label{sec:diffusion-experiment-summary}

The experiments support three conclusions.  First, relaxing privacy improves
finite-sample accuracy, with the dependence on $\alpha$ weakening as $n$
increases; at moderate and large privacy budgets, the OU-TV sample-size slope
is close to $-3/8$.  Second, the smoothness dependence of the OU-TV slopes
tracks $-(\gamma+1)/(2\gamma+3)$ across all five targets.  Third, the reverse
sampler reaches the finite-particle resolution floor well before the default
theoretical grid is exhausted, so private density estimation is the dominant
error source at the reported settings.

These findings are consistent with the theoretical rate decomposition. The target is a fixed finite
multiscale family, logarithmic factors and constants are visible at these
sample sizes, projected LMC is run for a finite number of steps, and the OU
score uses validated numerical interpolation rather than exact arithmetic.
Generated binned TV and Wasserstein distance are retained as supporting
metrics and are not used to redefine the population-TV theory.


\end{document}